\documentclass[10pt]{article}

% ---------- Document Identity (update these only) ----------
\newcommand{\DocStage}{}
\newcommand{\DocTitle}{Your Road to Tier One SOF}
\newcommand{\ProgramName}{182-Week Civilian-to-Selection Readiness Pipeline}
\newcommand{\ProgramSubtitle}{}

% ---------- Page ----------
\usepackage[legalpaper,left=0.5in,right=0.5in,top=0.6in,bottom=0.45in]{geometry}

% ---------- Typography ----------
\usepackage[T1]{fontenc}
\usepackage[utf8]{inputenc}
\usepackage{libertinus}
\usepackage{microtype}
\usepackage{parskip}
\usepackage{ragged2e}
\usepackage{textcomp}
\AtBeginDocument{\color{black}}
\AtBeginDocument{\pagecolor{white}}
\AtBeginDocument{\justifying}

% ---------- Landscape pages ----------
\usepackage{pdflscape}

% ---------- Color ----------
\usepackage[table]{xcolor}
\definecolor{StageBlue}{HTML}{1F4E79}
\definecolor{StageBlueDark}{HTML}{163A5A}
\definecolor{LightGray}{HTML}{F2F4F7}
\definecolor{MidGray}{HTML}{D9DEE5}
\definecolor{RuleGray}{HTML}{C7CED8}
\definecolor{HeaderInk}{HTML}{000000}
\definecolor{Ink}{HTML}{000000}

% ---------- Utility ----------
\usepackage{enumitem}
\sloppy
\setlength{\emergencystretch}{2em}

% ---------- Boxes / UI ----------
\usepackage[most]{tcolorbox}

\tcbset{
  charterTitle/.style={
    enhanced,
    colback=StageBlue!4,
    colframe=StageBlue,
    boxrule=1.25pt,
    arc=3pt,
    left=16pt,right=16pt,top=14pt,bottom=14pt,
    borderline north={3.2pt}{0pt}{StageBlue},
    borderline west={4pt}{0pt}{StageBlue},
    borderline south={1.25pt}{0pt}{StageBlue},
    drop shadow={black!25!white},
  },
  charterRule/.style={
    enhanced,
    colback=LightGray,
    colframe=RuleGray,
    boxrule=0.85pt,
    arc=3pt,
    left=12pt,right=12pt,top=10pt,bottom=10pt,
    borderline west={3pt}{0pt}{StageBlue},
  },
  charterBadge/.style={
    enhanced,
    colback=StageBlue,
    coltext=white,
    colframe=StageBlueDark,
    boxrule=0.6pt,
    arc=2pt,
    left=6pt,right=6pt,top=3pt,bottom=3pt,
  }
}
\newtcbox{\Badge}{nobeforeafter,charterBadge}

% ---------- Lists ----------
\setlist[itemize]{leftmargin=1.5em, itemsep=0.25em, topsep=0.25em}
\setlist[enumerate]{leftmargin=1.8em, itemsep=0.25em, topsep=0.25em}

% ---------- TOC ----------
\usepackage{tocloft}
\setcounter{tocdepth}{3}
\setlength{\cftbeforesecskip}{3pt}
\setlength{\cftbeforesubsecskip}{2pt}
\setlength{\cftbeforesubsubsecskip}{0pt}
\renewcommand{\cftsecfont}{\bfseries}
\renewcommand{\cftsecpagefont}{\bfseries}
\renewcommand{\cftsubsecfont}{\normalfont}
\renewcommand{\cftsubsecpagefont}{\normalfont}
\renewcommand{\cftsubsubsecfont}{\normalfont}
\renewcommand{\cftsubsubsecpagefont}{\normalfont}
\setlength{\cftsecindent}{0pt}
\setlength{\cftsubsecindent}{1.25em}
\setlength{\cftsubsubsecindent}{2.8em}
\setlength{\cftsecnumwidth}{2.1em}
\setlength{\cftsubsecnumwidth}{2.8em}
\setlength{\cftsubsubsecnumwidth}{3.6em}

% Remove the built-in TOC heading so only the custom heading is shown
\makeatletter
\renewcommand{\tableofcontents}{\@starttoc{toc}}
\makeatother

% ---------- Header/Footer: page number only ----------
\usepackage{fancyhdr}
\pagestyle{fancy}
\fancyhf{}
\renewcommand{\headrulewidth}{0pt}
\renewcommand{\footrulewidth}{0pt}
\fancyfoot[C]{\thepage}

% ---------- Section formatting ----------
\usepackage{titlesec}

\titleformat{\section}
  {\Large\bfseries\color{Ink}}
  {\thesection}{0.6em}{}
  [\vspace{0.25em}\textcolor{StageBlue}{\rule{\linewidth}{0.8pt}}]

\titleformat{\subsection}
  {\large\bfseries\color{Ink}}
  {\thesubsection}{0.6em}{}

\titleformat{\subsubsection}
  {\normalsize\bfseries\color{Ink}}
  {\thesubsubsection}{0.6em}{}

\titlespacing*{\section}{0pt}{0.95em}{0.35em}
\titlespacing*{\subsection}{0pt}{0.65em}{0.25em}
\titlespacing*{\subsubsection}{0pt}{0.55em}{0.2em}

% ---------- Table engine ----------
\usepackage{tabularray}
\UseTblrLibrary{booktabs}

% ---------- tabularray: GLOBAL table treatment ----------
\SetTblrInner{
  cells = {fg=black, bg=white, valign=t},
}

% ---------- Header helpers ----------
\newcommand{\Hdr}[1]{\shortstack[c]{\strut #1\strut}}

% ---------- Lines ----------
\newcommand{\TitleLine}{\textcolor{StageBlue}{\rule{0.98\linewidth}{0.95pt}}}
\newcommand{\SubTitleLine}{\textcolor{RuleGray}{\rule{0.98\linewidth}{0.65pt}}}

% ---------- Continued on next page ----------
\newcommand{\ContinuedOnNextPageInline}{%
  \par
  \vfill
  \noindent\textcolor{RuleGray}{\rule{\linewidth}{1pt}}\vspace{-2pt}
  \noindent\hfill\textit{Continued on next page}\par\vspace{-10pt}
  \noindent\textcolor{RuleGray}{\rule{\linewidth}{1pt}}
  \clearpage
}

% ---------- PDF hyperlinks ----------
\usepackage[hidelinks]{hyperref}
\hypersetup{
  colorlinks=false,
  pdfborder={0 0 0},
  linktoc=all,
  pdftitle={\DocTitle},
  pdfauthor={\ProgramName}
}

% =========================================================
\begin{document}

\subsection*{Phase 12 Card Header}
\phantomsection
\addcontentsline{toc}{subsection}{Phase 12 Card Header}

\begin{itemize}
\item Phase \textbf{ID}: Phase 12
\item Phase Name: Phase 12 — Final Pre-Selection Conditioning
\item Duration weeks: 12
\item Version: v1.0
\item Stage Alignment: Stage 5
\item Governing Status: Draft — Non-Deployable
\end{itemize}

\subsection*{1. Phase Mission}
\phantomsection


Execute a taper-aware verification phase that closes replication gaps and produces low-variance, gate-admissible evidence for \textbf{SELECTION-READY} classification without stacking high-cost tests or violating protected-window doctrine.

\subsection*{2. Phase Purpose}
\phantomsection


\begin{itemize}
\item Convert Phase 11 preparation into a controlled final verification block with minimal novelty and maximal comparability.
\item Close replication gaps for gate-critical metrics (2 of last 3 \textbf{VALID} for gate use per Stage 2 doctrine) without overtesting.
\item Establish a final “dominant” end gate posture (\textbf{RUN} + \textbf{CAL} or \textbf{RUCK} + \textbf{CAL}) rather than attempting maximal outcomes across all domains in one protected week.
\end{itemize}

\subsection*{3. Entry Criteria}
\phantomsection


\begin{itemize}
\item Entry requires admissible evidence from Phase 11 end gate consistent with Stage 1.F: no gate decisions based on \textbf{INVALID} tests or incomplete logs.
\item Required pre-entry posture:
  \begin{itemize}
  \item stable footwear/footcare status (no active blister; no unresolved gait-altering pain),
  \item environment/tool stability posture established (route/pool/tool identifiers stable per Stage 2.E),
  \item adherence posture stable (decision-window compliance per Stage 1.F).
  \end{itemize}
\item If evidence is incomplete or safety posture is unstable, default posture is \textbf{HOLD} and reslot per Stage 3.
\end{itemize}

\subsection*{4. Operating Constraints}
\phantomsection


\begin{itemize}
\item Six sessions per week is mandatory.
\item Required session elements per Stage 1.F in order:
  \begin{itemize}
  \item Safety self-check first
  \item Warm-up
  \item Mobility
  \item Main work
  \item Cardio
  \item Cooldown
  \item Recovery notes and any required post-session notes per Stage 1.F
  \end{itemize}
\item Cardio minimum standard: minimum 30 minutes continuous per session.
\item 10,000 steps per day global requirement.
\item Validity doctrine: gate decisions rely on admissible Valid evidence only; \textbf{INVALID} tests provide no gate credit; sessions must meet validity posture per Stage 1.F for gate eligibility.
\end{itemize}

\subsection*{5. Primary Adaptations and Physiological Focus}
\phantomsection


\begin{itemize}
\item Verification quality: reduce variance and produce repeatable \textbf{VALID} outcomes under stable environment/tool conditions.
\item Specificity: emphasize selection-relevant tests without introducing new instruments.
\item Taper-aware freshness management: protect high-cost tests from fatigue contamination.
\item Durability stability: preserve foot integrity and gait mechanics as gate constraints.
\item Replication closure: schedule enough protected windows to satisfy replication requirements without compression.
\end{itemize}

\subsection*{6. Primary Modalities}
\phantomsection


\begin{itemize}
\item Emphasized domains: \textbf{RUN}, \textbf{CAL}, \textbf{RUCK}, \textbf{SWIM}, \textbf{DURABILITY}, \textbf{RECOVERY}, \textbf{BODYCOMP}.
\item Supporting domains (only if already established and non-disruptive): \textbf{STR} (submax proxy), \textbf{AER} (conditioning domain evidence; not interchangeable with \textbf{RUN}).
\item De-emphasized/prohibited in Phase 12:
  \begin{itemize}
  \item any new test constructs,
  \item any new domain tags,
  \item any change to Stage 2 protocol rules,
  \item any underwater work unless governance prerequisites are satisfied (and it must not compromise dominant gate evidence).
  \end{itemize}
\end{itemize}

\clearpage
\begin{landscape}

\subsection*{7. Weekly Structure}
\phantomsection


\begingroup
\scriptsize
\setlength{\tabcolsep}{2pt}
\renewcommand{\arraystretch}{1.12}

\begin{longtblr}{
  width=\linewidth,
  rowhead=1,
  colspec = {
    Q[l,wd=0.70in]
    X[l,2.05]
    X[l,1.25]
    X[l,2.95]
    X[l,2.90]
  },
  hlines,
  vlines,
  cells = {hsep=6pt, vsep=6pt, halign=l, valign=t},
  cell{1-Z}{1-5} = {fg=black},
  row{1} = {bg=StageBlue!15, fg=black, font=\bfseries, halign=l, valign=t},
  row{odd} = {bg=white},
  row{even} = {bg=LightGray},
}
\Hdr{Session \#} &
\Hdr{Primary Focus} &
\Hdr{Main Work Domains} &
\Hdr{Cardio Modality/Intent} &
\Hdr{Notes/Risk Controls} \\

1 &
\textbf{RUN} readiness and mechanics stability &
\textbf{RUN}; \textbf{DURABILITY} &
Modality: step-based walk or treadmill walk; Minutes: 30–60; RPE plus Talk Test: RPE 3–5, full to short sentences; Continuity: continuous planned &
No novelty in footwear/route/tool. Any gait drift constrains upcoming test exposure. \\

2 &
\textbf{CAL} maintenance and evidence integrity &
\textbf{CAL}; \textbf{RECOVERY} &
Modality: low-impact steady; Minutes: 30–45; RPE plus Talk Test: RPE 2–4, full sentences; Continuity: continuous planned &
\textbf{CAL} artifacts required when \textbf{CAL} tests are executed in protected windows. Avoid shoulder flare-ups. \\

3 &
\textbf{RUCK} tolerance maintenance and foot integrity &
\textbf{RUCK}; \textbf{DURABILITY} &
Modality: step-based walk or treadmill walk; Minutes: 30–60; RPE plus Talk Test: RPE 3–5, short sentences permitted but controlled; Continuity: continuous planned &
No load escalation implied here. Foot hotspot/blister posture is decisive. \\

4 &
\textbf{SWIM} readiness or controlled substitute posture &
\textbf{SWIM}; \textbf{RECOVERY} &
Modality: low-impact steady; Minutes: 30–45; RPE plus Talk Test: RPE 2–4, full sentences; Continuity: continuous planned &
Pool access and unit verification are prerequisites for \textbf{SWIM} testing. If pool not feasible, reslot \textbf{SWIM} tests; no proxy equivalency. \\

5 &
Recovery stabilization day &
\textbf{RECOVERY}; \textbf{DURABILITY} &
Modality: lowest-risk continuous; Minutes: 30–45; RPE plus Talk Test: RPE 2–3, full sentences; Continuity: continuous planned &
Protect freshness for protected windows; no “make-up” intensity. \\

6 &
Consolidation and pre-window stabilization &
\textbf{RECOVERY}; \textbf{DURABILITY} &
Modality: low-impact steady; Minutes: 30–45; RPE plus Talk Test: RPE 2–4, full sentences; Continuity: continuous planned &
Enforce no-novelty posture before Deload+Test and Transition weeks. \\

\end{longtblr}

\endgroup

\subsection*{8. Progression Rules}
\phantomsection


\begin{itemize}
\item Phase 12 progression is verification-driven, not expansion-driven:
  \begin{itemize}
  \item minimize novelty and variance,
  \item prioritize stable execution and reproducibility over “bigger numbers.”
  \end{itemize}
\item One-change rule in the lead-in to protected windows:
  \begin{itemize}
  \item do not change route, treadmill, footwear, pool unit/length, timing tool, or test setup inside a decision window.
  \end{itemize}
\item Reslot posture:
  \begin{itemize}
  \item if a high-cost test is \textbf{INVALID} or unsafe, reslot to the next protected window; do not compress by stacking.
  \end{itemize}
\item Durability constraints dominate:
  \begin{itemize}
  \item gait-altering pain flag drives \textbf{HOLD} posture for locomotion progression decisions,
  \item foot hotspot/blister drift forces conservative substitution and may reslot ruck/run tests.
  \end{itemize}
\end{itemize}

\newpage

\subsection*{9. Deload and Test Placement}
\phantomsection


Phase 12 is taper-aware per Stage 3 integrated calendar posture:

\begin{itemize}
\item Week 1 baseline (low-risk baseline touchpoint; not a maximal gate by default).
\item Week 5 Deload+Test (mid-phase decision touchpoint).
\item Week 10 Deload+Test (verification window; used to close replication gaps).
\item Week 11 Transition (stabilization and taper-aware posture; no new major tests).
\item Week 12 Deload+Test (end-of-phase exit gate window).
\end{itemize}

\end{landscape}

\clearpage
\begin{landscape}

\subsection*{10. Test Battery}
\phantomsection


\begingroup
\scriptsize
\setlength{\tabcolsep}{2pt}
\renewcommand{\arraystretch}{1.12}

\begin{longtblr}{
  width=\linewidth,
  rowhead=1,
  colspec = {
    X[l,1.55]
    X[l,3.05]
    Q[l,wd=0.95in]
    X[l,3.35]
    Q[l,wd=1.35in]
  },
  hlines,
  vlines,
  cells = {hsep=6pt, vsep=6pt, halign=l, valign=t},
  cell{1-Z}{1-5} = {fg=black},
  row{1} = {bg=StageBlue!15, fg=black, font=\bfseries, halign=l, valign=t},
  row{odd} = {bg=white},
  row{even} = {bg=LightGray},
}
\Hdr{Test Timing} &
\Hdr{Test \textbf{ID}/Name} &
\Hdr{Domain} &
\Hdr{Validity Conditions} &
\Hdr{Pass Threshold Stage 2 Tier} \\

Week 5 Deload+Test &
\textbf{C9} / \textbf{MET-2B-015} 2-Mile Run Time Trial &
\textbf{RUN} &
Stage 2 protocol conditions; stable Route \textbf{ID} or Machine \textbf{ID}; warm-up logged; no stop-rule violations; complete Stage 2.E fields &
\textbf{INTERMEDIATE} \\

Week 5 Deload+Test &
\textbf{C4} / \textbf{MET-2B-010} Push-Ups — 2-Minute Timed, Strict &
\textbf{CAL} &
Stage 2 protocol; \textbf{CAL} artifact required for gate-validity; warm-up logged; rep standard auditable &
\textbf{INTERMEDIATE} \\

Week 5 Deload+Test &
\textbf{C6} / \textbf{MET-2B-012} Pull-Ups — Strict Dead-Hang &
\textbf{CAL} &
Stage 2 protocol; \textbf{CAL} artifact required; pronated standard; ROM auditable &
\textbf{INTERMEDIATE} \\

Week 5 Deload+Test &
\textbf{C7} / \textbf{MET-2B-013} Plank — Forearm &
\textbf{CAL} &
Stage 2 protocol; \textbf{CAL} artifact required; timer auditable; termination on form failure &
\textbf{INTERMEDIATE} \\

Week 10 Deload+Test &
\textbf{C12} / \textbf{MET-2B-019} 6-Mile Ruck Time Trial — 35 lb dry &
\textbf{RUCK} &
Stage 2 protocol; dry load verified; water logged separately; Route \textbf{ID} or Machine \textbf{ID}; footwear logged; no pauses; no stop-rule violations &
\textbf{INTERMEDIATE} \\

Week 10 Deload+Test &
\textbf{C13} / \textbf{MET-2B-022} Swim 500 yd Time — No Fins OR \textbf{C13} / \textbf{MET-2B-023} Swim 500 m Time — No Fins &
\textbf{SWIM} &
Unit locked; pool verification logged; no fins; allowable strokes; rest governance per protocol; no stop-rule violations &
\textbf{INTERMEDIATE} \\

Week 12 Deload+Test &
\textbf{C11} / \textbf{MET-2B-017} 5-Mile Run Time Trial &
\textbf{RUN} &
Stage 2 protocol; stable Route \textbf{ID} or Machine \textbf{ID}; warm-up logged; no pauses; no stop-rule violations; environment logged &
\textbf{SELECTION-READY} \\

Week 12 Deload+Test &
\textbf{C4} / \textbf{MET-2B-010} Push-Ups — 2-Minute Timed, Strict &
\textbf{CAL} &
Stage 2 protocol; \textbf{CAL} artifact required; warm-up logged &
\textbf{SELECTION-READY} \\

Week 12 Deload+Test &
\textbf{C5} / \textbf{MET-2B-011} Sit-Ups — 2-Minute Timed, Standard &
\textbf{CAL} &
Stage 2 protocol; \textbf{CAL} artifact required; warm-up logged &
\textbf{SELECTION-READY} \\

Week 12 Deload+Test &
\textbf{C6} / \textbf{MET-2B-012} Pull-Ups — Strict Dead-Hang &
\textbf{CAL} &
Stage 2 protocol; \textbf{CAL} artifact required; pronated standard &
\textbf{SELECTION-READY} \\

Week 12 Deload+Test &
\textbf{C7} / \textbf{MET-2B-013} Plank — Forearm &
\textbf{CAL} &
Stage 2 protocol; \textbf{CAL} artifact required; timer auditable &
\textbf{SELECTION-READY} \\

Week 12 Deload+Test (conditional) &
\textbf{C14} / \textbf{MET-2B-025} Underwater Exposure Admissibility Flag; \textbf{C14} / \textbf{MET-2B-026} Underwater 25 m Time — One Breath &
\textbf{SWIM} &
Only if governance prerequisites are satisfied (certified lifeguard + dedicated safety spotter; facility permission); otherwise not scheduled; no proxy equivalency &
\textbf{SELECTION-READY} \\

\end{longtblr}

\endgroup

\end{landscape}
\clearpage

\subsection*{11. Exit Standards}
\phantomsection


\begin{itemize}
\item Gate outcome to advance: \textbf{ADVANCE} only when Phase 12 end gate evidence is admissible and meets the \textbf{SELECTION-READY} tier posture for the listed decisive tests.
\item End gate is Week 12 Deload+Test. If Week 12 evidence is \textbf{INVALID} or missed, default posture is \textbf{HOLD} and reslot per Stage 3.
\item Intended tier posture alignment for Phase 12 exit: \textbf{SELECTION-READY}.
\item Replication posture: tier classification for progression decisions relies on replication requirements (2 of last 3 \textbf{VALID} in the decision window); where replication is incomplete, posture defaults \textbf{HOLD} for that metric until replication exists.
\end{itemize}

\subsection*{12. Risk Controls and Stop Rules}
\phantomsection


\begin{itemize}
\item Stage 0.5 and Stage 1.F stop posture controls; not rewritten here.
\item Phase 12-specific high-risk controls:
  \begin{itemize}
  \item \textbf{RUN} and \textbf{RUCK} maximal events must not be stacked in a way that violates Stage 3 anti-stacking/spacing intent; prioritize dominant gate evidence and reslot secondary events.
  \item Foot integrity is gate-critical: hotspots/blisters and gait drift constrain ruck/run test admissibility; do not “push through.”
  \item \textbf{SWIM} validity depends on unit-locked pool verification; pool disruption triggers reslot rather than compromised attempts.
  \item Underwater remains governance-gated and is conditional only; if governance is not guaranteed, it is not scheduled and is inadmissible without governance.
  \end{itemize}
\end{itemize}

\clearpage
\begin{landscape}

\subsection*{13. Required Capabilities and Dependencies}
\phantomsection


\begingroup
\scriptsize
\setlength{\tabcolsep}{2pt}
\renewcommand{\arraystretch}{1.12}

\begin{longtblr}{
  width=\linewidth,
  rowhead=1,
  colspec = {
    X[l,1.85]
    X[l,3.15]
    Q[l,wd=1.25in]
    Q[l,wd=1.25in]
    X[l,3.30]
  },
  hlines,
  vlines,
  cells = {hsep=6pt, vsep=6pt, halign=l, valign=t},
  cell{1-Z}{1-5} = {fg=black},
  row{1} = {bg=StageBlue!15, fg=black, font=\bfseries, halign=l, valign=t},
  row{odd} = {bg=white},
  row{even} = {bg=LightGray},
}
\Hdr{Dependency \textbf{ID}} &
\Hdr{Required Capability Category and Tier} &
\Hdr{Due-By week-in-phase} &
\Hdr{Owner} &
\Hdr{Fail Action} \\

\textbf{DEP-1C-013} &
7.07 Load-Bearing Rucking and Carry Systems — Standard &
Week 4 &
Equipment Lead &
\textbf{HOLD} the Week 10 ruck gate; reslot to next protected window; no proxy equivalency; preserve Cardio and durability controls. \\

\textbf{DEP-1C-007} &
7.11 Footwear Systems and Foot Care — Standard &
Week 1 &
Trainee &
If hotspots/blisters present or footwear unstable: downshift locomotion; \textbf{HOLD} ruck/run gate attempts until stable; reslot. \\

\textbf{DEP-1C-006} &
7.18 Testing Timing Measurement and Course-Setup Tools — Standard &
Week 1 &
Coach Lead &
If tools/route/pool verification cannot be logged: treat tests \textbf{INVALID}; reslot; do not claim tier attainment. \\

\textbf{DEP-1C-011} &
7.17 Aquatic Training and Water Safety Equipment — Standard &
Week 5 &
Equipment Lead &
If pool unit cannot be verified or access unstable: reslot \textbf{SWIM} \textbf{TT}; do not substitute with non-swim evidence for \textbf{SWIM} gate. \\

\textbf{DEP-1C-017} &
7.12 Recovery Mobility and Tissue-Care Implements — Standard &
Week 1 &
Trainee &
If recovery drift rises (sleep collapse, soreness trend): convert to Minimum Viable Sessions; protect protected windows; reslot maximal tests if validity would be compromised. \\

\textbf{DEP-1C-012} &
7.01 Strength and Resistance Training Equipment — Minimal &
Week 5 &
Equipment Lead &
If \textbf{STR} proxy cannot be executed safely/loggably: omit \textbf{STR} proxy in Phase 12; do not replace locomotion gates; record as not executed (supporting only). \\

\textbf{DEP-1C-015} &
7.09 Health Monitoring Diagnostics and Biometrics — Minimal &
Week 1 &
Coach Lead &
If logs are incomplete: evidence is inadmissible; \textbf{HOLD} gate decisions; restore logging discipline before attempting verification windows. \\

\textbf{DEP-1C-022} &
7.19 Protective and Assistive Equipment — Minimal &
Week 5 &
Equipment Lead &
If foot/skin protection capability is missing: increase risk of breakdown; downshift and reslot as needed; do not force ruck mileage. \\

\textbf{DEP-1C-016} &
7.10 Logistics Maintenance Storage and Sustainment Equipment — Minimal &
Week 1 &
Trainee &
If equipment readiness drifts due to logistics failures: treat as capability constraint; reslot tests rather than improvising with unverified setups. \\

\end{longtblr}

\endgroup

\subsection*{14. Validity Requirements}
\phantomsection


\begin{itemize}
\item Admissible evidence for gates requires Stage 1.F validity tagging discipline and Stage 2 validity conditions for each protocol card.
\item \textbf{INVALID} tests provide no gate credit and require reslot per Stage 3; no advancement claims are permitted off invalid evidence.
\item \textbf{CAL} tests require compliant artifact evidence for \textbf{VALID} gate use; missing artifacts render the test \textbf{INVALID} for gate use.
\item Environment/tool stability per Stage 2.E is mandatory for comparability: route/tool/pool verification and consistent logging fields are required.
\end{itemize}

\end{landscape}

\subsection*{15. Change Control Notes}
\phantomsection


\begin{itemize}
\item Allowed without patch: reslot a missed/invalid test to the next protected window; reorder tests within a protected week while preserving Stage 3 anti-stacking/spacing intent; downgrade a protected week to deload-only/check-in-only if safety/validity is compromised.
\item Requires patch: phase duration change; adding protected windows; changing gate test set; changing dependencies that affect gate eligibility.
\item Reslot posture follows Stage 3.E: no stacking make-ups; reslot forward; default \textbf{HOLD} if end gate cannot be executed validly.
\end{itemize}

\subsection*{16. Compliance Checklist}
\phantomsection


\begin{itemize}
\item Pass if: all tests in the battery are Stage 2-defined (C\# and \textbf{MET-2B}-###), and protected windows are specified consistent with Stage 3.
\item Pass if: Cardio and steps constraints are stated in body text and weekly structure Cardio cells include modality/minutes/RPE+talk test/continuity.
\item Pass if: dependencies are category-level and use Stage 7 IDs; 7.02 is not required.
\item Fail (Draft — Non-Deployable) if: dependency IDs are placeholders affecting gating (present here), or any gate-relevant field contains “use Stage 2 \textbf{ID}” or \textbf{TBD} language.
\item Pass if: risk controls reference Stage 0.5/Stage 1.F authority and encode Stage 3 anti-stacking/spacing intent without programming.
\end{itemize}

\end{document}